\documentclass[12pt]{article}

% Layout.
\usepackage[top=1in, bottom=0.75in, left=1in, right=1in, headheight=1in, headsep=6pt]{geometry}

% Fonts.
\usepackage{mathptmx}
\usepackage[scaled=0.86]{helvet}
\renewcommand{\emph}[1]{\textsf{\textbf{#1}}}

% TiKZ.
\usepackage{tikz, pgfplots}
\usetikzlibrary{calc}

% Misc packages.
\usepackage{amsmath,amssymb,latexsym}
\usepackage{graphicx}
\usepackage{array}
\usepackage{xcolor}
\usepackage{multicol}
\usepackage{fancyhdr}
\pagestyle{fancy}


% Misc.
\renewcommand{\d}{\displaystyle}
\newcommand{\ds}{\displaystyle}
\newcommand{\ul}[1]{\underline{#1}}
\def\bc{\begin{center}}
\def\ec{\end{center}}
\def\be{\begin{enumerate}}
\def\ee{\end{enumerate}}

\usepackage{array}
\newcolumntype{L}[1]{>{\raggedright\let\newline\\\arraybackslash\hspace{0pt}}m{#1}}
\newcolumntype{C}[1]{>{\centering\let\newline\\\arraybackslash\hspace{0pt}}m{#1}}
\newcolumntype{R}[1]{>{\raggedleft\let\newline\\\arraybackslash\hspace{0pt}}m{#1}}
\newcommand{\red}[1]{\textcolor{red}{#1}}

\newcommand{\ans}[1][1in]{\rule{#1}{.5pt}}


\lhead{Math F113X: Quiz 7}
\rhead{14 November, 2025}

\begin{document}

\vspace{1cm}

\noindent \textbf{Name:} \hrulefill \quad  score:\ans[1cm] \ / 10 \\


There are 10 points possible on this quiz. No aids (book, notes, etc.)
are permitted. You may use a non-programmable calculator.  {\bf Show all work and supporting calculations for full credit. Explain how you get your answers.}

\vspace{0.2cm}

A tabula recta is attached.


\be

\item (2 points)  \textbf{Encrypt} the message SCHEDULING using an alphabetic Caesar cipher with shift 4 (A to E).
%NYXKWNIFD

\vfill

\item (2 points)  \textbf{Decrypt} the message NYHWO AOLVYF using an alphabetic Caesar cipher with shift 7 (A to H).
%GRAPH THEORY

\vfill

\item (2 points)  \textbf{Encrypt} the message 14 NOV 2025 using the substitution table below.\\

{
\begin{tabular}{|c|c|c|c|c|c|c|c|c|c|c|c|c|c|c|c|c|c|c|}
\hline
%&0&1&2&3&4&5&6&7&8&9&10&11&12&13&14&15&16&17\\ \hline \hline
original& A&B&C&D&E&F&G&H&I&J&K&L&M&N&O&P&Q&R\\  \hline
maps to&9&8&7&6&5&4&3&2&1&0&D&C&B&A&H&G&F&E\\
\hline
\end{tabular}

\begin{tabular}{|c|c|c|c|c|c|c|c|c|c|c|c|c|c|c|c|c|c|c|}
 \hline
%&18&19&20&21&22&23&24&25&26&27&28&29&30&31&32&33&34&35\\ \hline \hline
original &S&T&U&V&W&X&Y&Z&0&1&2&3&4&5&6&7&8&9\\ \hline
maps to&L&K&J&I&P&O&N&M&T&S&Q&R&V&U&Y&X&W&Z \\
%[12pt]
\hline
\end{tabular}
}

\vfill
\newpage
\item (2 points)  \textbf{Encrypt} the message NEED MORE COFFEE using a tabular transposition cipher with rows of length 5.
  

\vfill

\item (3 points)   \textbf{Decrypt} the message TSHHH FEWAO FXNRI YKTSZ using a tabular transposition cipher with rows of length 5.
%THANKS FOR THE FISH

\vfill
\item (1 point)  Describe at least one aspect of a substitution cipher that makes it vulnerable to being hacked (that is, deciphered by an unwanted observer).
\vfill

\ee

\newpage

\makeatletter
\newcommand{\Letter}[1]{\@Alph{#1}}
\makeatother

\def\r{.6}
\begin{tikzpicture}
\draw[step=\r] (0,0) grid (28*\r,28*\r);
%\foreach \i in {1,2,...,26}{\draw (0,0) -- (\i*\r, 1);}
\foreach \i in {1,2,...,26}{\path let \n1 = {int(\i-1)} in (\i*\r+1/2*\r+1*\r, 28*\r-1/2*\r) node {\n1};}
\foreach \i in {1,2,...,26}{\path (\i*\r+1/2*\r +1*\r, 27*\r-1/2*\r) node {\textbf{\Letter{\i}}};}%
\foreach \i in {1,2,...,26}{\path (1/2*\r+1*\r, 26*\r-\i*\r+1/2*\r) node {\textbf{\Letter{\i}}};}%
\foreach \i in {1,2,...,26}{\path let \n1 = {int(\i-1)} in  (1/2*\r, 26*\r-\i*\r+1/2*\r) node {\n1};}%
\draw[line width = .75mm] (0,26*\r) -- (28*\r, 26*\r);
\draw[line width = .75mm] (\r+1*\r,0) -- (\r+1*\r, 28*\r);
\foreach \i in {1,2,...,26}
	\foreach \j in {1,2,...,26}{
		{\path let \n1 = {int(mod((\j-1)+(\i-1), 26)+1} in (\i*\r+1/2*\r+1*\r, 27*\r-\j*\r - 1/2*\r) node {\Letter{\n1}};
		}}
\end{tikzpicture}


\end{document}
%%% Local Variables:
%%% mode: LaTeX
%%% TeX-master: t
%%% End:
